\documentclass[a4paper,10pt]{article}
\usepackage[UTF8]{ctex}
\usepackage{geometry}
\usepackage{tikz}
\usetikzlibrary{shapes.geometric, arrows.meta, positioning, calc, trees}
\usepackage{longtable}
\usepackage{array}
\usepackage{fancyhdr}
\usepackage{xcolor}
\usepackage{amsmath}
\usepackage{amssymb}
\usepackage{hyperref}
\geometry{left=2.5cm, right=2.5cm, top=3cm, bottom=2.5cm}
\pagestyle{fancy}
\fancyhf{}
\fancyhead[C]{\small clampControlInput 函数设计文档}
\fancyfoot[C]{\thepage}
\setcounter{secnumdepth}{3}
\setcounter{tocdepth}{3}
\tikzset{
    startstop/.style={ellipse, draw, fill=red!10, text width=8em, minimum height=0.9cm, align=center},
    process/.style={rectangle, draw, fill=blue!5, text width=11em, minimum height=1.1cm, align=center},
    decision/.style={diamond, draw, fill=green!10, text width=7em, minimum height=1cm, align=center},
    arrow/.style={thick,->,>=Stealth}
}
\begin{document}
\begin{titlepage}
    \centering
    {\Huge\textbf{clampControlInput（控制输入双轴饱和限幅）函数设计}\par}
    \vspace{1.5cm}
    \renewcommand{\arraystretch}{1.5}
    \begin{tabular}{|l|l|l|l|}
        \hline
        文件密级 & 内部公开 & 编号 & CBF-DOC-CLAMPCONTROLINPUT \\ \hline
        版本 & V1.0 & 内容 & clampControlInput（控制输入双轴饱和限幅）函数设计 \\ \hline
        设计 & 许庆 & 联系方式 & 18612426814 / qingxu@tsinghua.edu.cn \\ \hline
        审核 & - & 审批 & - \\ \hline
        发布日期 & 2026 年 6 月 26 日 & & \\ \hline
    \end{tabular}
    \vspace{1.5cm}
    \section*{文件更改记录}
    \begin{longtable}{|c|c|p{4cm}|c|c|c|c|}
        \hline
        日期 & 版本号 & 修订说明 & 辅助AI & 修订 & 审核 & 审批 \\ \hline
        2026-6-26 & V1.0 & 初次设计 & Claude-Opus-4.7 & 许庆 & - & - \\ \hline
    \end{longtable}
\end{titlepage}
\tableofcontents
\newpage
\section{公有函数设计}
\subsection{clampControlInput（控制输入双轴饱和限幅）函数设计}
\subsubsection{函数需求}
\par 本函数实现 \textbf{控制输入双轴饱和限幅}。其在 CBF 自动驾驶安全控制流水线中承担如下职责：对 (a, $\textbackslash{}delta$\_f) 做 box 限幅，确保不越界。
\par 头文件路径：\texttt{include/cpp/clampControlInput/clampControlInput.hpp}；源文件路径：\texttt{src/cpp/clampControlInput/clampControlInput.cpp}。
\subsubsection{算法设计}
\par 本函数核心数学表达式如下：
\begin{equation}
a^* = \max(a_{min}, \min(a_{max}, a)),\quad \delta^*_f = \max(\delta_{min}, \min(\delta_{max}, \delta_f))
\end{equation}
\par 实现步骤：
\begin{enumerate}
  \item \textbf{clamp\_a}: 对加速度 std::clamp
  \item \textbf{clamp\_delta}: 对转角 std::clamp
\end{enumerate}
\subsubsection{程序流程图}
\begin{figure}[h]
    \centering
    \begin{tikzpicture}[node distance=1.6cm]
    \node (start) [startstop] {开始};
    \node (p0) [process, below of=start, yshift=-0.3cm] {对加速度 std::clamp};
    \node (p1) [process, below of=p0, yshift=-0.3cm] {对转角 std::clamp};
    \node (end) [startstop, below of=p1, yshift=-0.3cm] {返回 / 写回输出};
    \draw [arrow] (start) -- (p0);
    \draw [arrow] (p0) -- (p1);
    \draw [arrow] (p1) -- (end);

    \end{tikzpicture}
    \caption{clampControlInput 函数程序流程图}
    \label{fig:flow_clampControlInput}
\end{figure}
\subsubsection{函数接口}
\par 输入参数列表：
\begin{itemize}
  \item \texttt{aRaw} (double) -- 限幅前加速度
  \item \texttt{deltaFRaw} (double) -- 限幅前前轮转角
  \item \texttt{aMin/aMax} (double) -- 加速度上下界
  \item \texttt{deltaFMin/deltaFMax} (double) -- 转角上下界
\end{itemize}
\par 输出参数列表：
\begin{itemize}
  \item \texttt{aFinal} (double*) -- 限幅后加速度
  \item \texttt{deltaFFinal} (double*) -- 限幅后转角
\end{itemize}
\subsubsection{输入输出模块图}
\begin{figure}[h]
    \centering
    \begin{tikzpicture}[>=Stealth]
        \node[draw, rectangle, minimum width=4.5cm, minimum height=2cm, fill=blue!5] (func) at (0,0) {\large{\textbf{clampControlInput}}};
        \draw[->, thick] (-3.5,0.4) -- node[above]{I (输入)} (-2.3,0.4);
        \draw[->, thick] (2.3,0.4) -- node[above]{O (输出)} (3.5,0.4);
        \draw[->, thick] (0,-1.0) -- node[right]{f (返回)} (0,-2.0);
    \end{tikzpicture}
    \caption{clampControlInput 函数输入输出模块图}
    \label{fig:io_clampControlInput}
\end{figure}
\subsubsection{函数诸元表}
\begin{table}[h]
    \centering
    \caption{clampControlInput 函数诸元表}
    \label{tab:meta_clampControlInput}
    \begin{tabular}{|c|c|c|c|c|}
        \hline
        中文名 & 英文名 & 性质 & 级别 & 入库 \\ \hline
        控制输入双轴饱和限幅 & clampControlInput & 元件 & 应用层 & 是 \\ \hline
    \end{tabular}
\end{table}
\subsubsection{函数架构}
\par 本函数为\textbf{元件}（叶子节点），无下级函数依赖。

\subsubsection{测试用例表}
\noindent\textbf{正常场景测试用例如表~\ref{tab:normal_clampControlInput} 所示：}
\begin{table}[h]
\centering
\caption{clampControlInput 正常场景测试用例}
\label{tab:normal_clampControlInput}
\begin{tabular}{|c|p{3.5cm}|p{3.5cm}|p{3.5cm}|}
\hline
ID & 描述 & 输入 & 期望输出 \\ \hline
  TC1 & 正常用例 1 & 见 cases.json & 见 cases.json \\ \hline
  TC2 & 正常用例 2 & 见 cases.json & 见 cases.json \\ \hline
  TC3 & 正常用例 3 & 见 cases.json & 见 cases.json \\ \hline
  TC4 & 正常用例 4 & 见 cases.json & 见 cases.json \\ \hline
  TC5 & 正常用例 5 & 见 cases.json & 见 cases.json \\ \hline
\end{tabular}
\end{table}

\noindent\textbf{异常场景测试用例如表~\ref{tab:abnormal_clampControlInput} 所示：}
\begin{table}[h]
\centering
\caption{clampControlInput 异常场景测试用例}
\label{tab:abnormal_clampControlInput}
\begin{tabular}{|c|p{3.5cm}|p{3.5cm}|p{3.5cm}|}
\hline
ID & 描述 & 输入 & 期望输出 \\ \hline
  EC1 & NaN 输入 & 任一标量=NaN & 输出含 NaN \\ \hline
  EC2 & Inf 输入 & 任一标量=Inf & 输出含 Inf \\ \hline
  EC3 & 数学边界 & 极限值 & 结果有界或返回 false \\ \hline
  EC4 & QP 不可行 & 约束相互冲突 & 返回 feasible=false \\ \hline
\end{tabular}
\end{table}

\noindent\textbf{符号说明：}
\begin{itemize}
  \item dx, dy: ego-frame 相对位置 (m)
  \item v\_rx, v\_ry: ego-frame 相对速度 (m/s)
  \item r: 安全圆半径 (m)
  \item B: Barrier 函数值
\end{itemize}
\end{document}
